%%% =====================================================================
%%%  Wave CFG2026 — Agent 4
%%%  Transverse E_2-shadow of the d=3 construction at chain level
%%%
%%%  Focal axis: At d=3 the chiral output A = Phi_3^{(Sigma_2,C)}(C)
%%%  is an E_1-algebra along the curve C.  The E_2-structure lives
%%%  on Z^der(Rep^{E_1}(A)), the derived Drinfeld centre of its
%%%  representation category, and is forced by the Dunn-Lurie
%%%  tensor-decomposition E_3 = E_1 x E_2 of the ambient native
%%%  E_3-holomorphic-factorisation algebra Phi^FA_3(C).
%%%
%%%  Chain-level derivation:
%%%    (i)   Dunn-Lurie additivity E_3 = E_1 (x) E_2  (Lurie HA 5.1.2.2).
%%%    (ii)  Derived Drinfeld centre: Z^der(Rep^{E_1}(A))
%%%          = End_{A (x) A^op}(A) in the dg-model.
%%%    (iii) Explicit R-matrix from CY_3 input via
%%%          Maulik-Okounkov stable envelopes, recognised as the
%%%          residue of the UV positive-half gluing cocycle
%%%          phi^+_UV across the equivariant-chamber wall.
%%%    (iv)  The braided monoidal structure on Z^der(Rep^{E_1}(A))
%%%          verifies the hexagon axioms chain-level and identifies
%%%          with Rep of the Drinfeld double D(Y^+(gl_1-hat)) for
%%%          C^3, and with Rep of Y_{osp}(4|20) in the K3 super-case.
%%%
%%%  Written with Etingof / Kazhdan / Drinfeld / Bezrukavnikov voices
%%%  and read side-by-side with Costello-Francis-Gwilliam 2026.
%%% =====================================================================

\section*{Agent 4.  The transverse E$_2$-shadow at d $=$ 3}

The single-sentence statement.  At $d = 3$ the CY-to-chiral functor
produces, through its first-stage $E_3$-holomorphic-factorisation
algebra $\PhiFA_3(\cC)$ on the CY$_3$ variety, \emph{two} algebraic
shadows after the two-stage factorisation.  One is $E_1$-chiral along
the specialised curve $C$; it is the familiar
$A = \Phi_3^{(\Sigma_2, C)}(\cC)$ of Chapter~\ref{ch:cy-to-chiral}.
The \emph{other}, canonically and Dunn--Lurie-forced, is an
$E_2$-algebra that does not live on $A$: it lives on the derived
Drinfeld centre $Z^{\mathrm{der}}(\Rep^{E_1}(A))$ of the representation
category of $A$.  This is the transverse $E_2$-shadow.  It is the
origin of the quantum-group $R$-matrix in Vol~III.

\subsection*{Ghost theorems (the five attack vectors)}

The five adversarial attacks each generate a ghost theorem, and the
surviving stable core is a chain-level statement we derive below.

\begin{enumerate}[label=\textup{(G\arabic*)}]
 \item \emph{Operadic-level ghost.} The $E_3$-hFA $\PhiFA_3(\cC)$ on
 the CY$_3$ variety admits, up to contractible choice, a Dunn--Lurie
 decomposition $E_3 \simeq E_1 \otimes_{E_0} E_2$, and hence two
 commuting algebraic refinements $(A, Z^{\mathrm{der}}(\Rep^{E_1}(A)))$.
 \item \emph{Representation-centre ghost.} The $E_2$-factor acts not on
 $A$ but on the derived Drinfeld centre of $\Rep^{E_1}(A)$:
 equivalently, on the Hochschild cochain complex of $A$ as a dg-algebra
 in the $E_1$-factorisation category along $C$.
 \item \emph{$R$-matrix ghost.} The $E_2$-braiding on
 $Z^{\mathrm{der}}(\Rep^{E_1}(A))$ is \emph{the} Maulik--Okounkov
 $R$-matrix, read as the residue of the UV positive-half gluing cocycle
 $\phi^+_{\mathrm{UV}}(u)$ across the equivariant-chamber wall at
 $u_\star$, and the MO axiom is the cocycle condition for
 $\phi^+_{\mathrm{UV}}$.
 \item \emph{Quasi-triangular-Hopf ghost.} For $\cC = \Perf(\C^3)$, the
 transverse $E_2$-shadow is canonically the Drinfeld double
 $D(Y^+(\widehat{\mathfrak{gl}}_1)) = Y(\widehat{\mathfrak{gl}}_1)$;
 for $\cC = D^b(\Coh(K3 \times E))$ in the reduced branch it is
 $Y_{\mathfrak{osp}}(4|20)$; for toric CY$_3$ more general it is the
 quantum-toroidal double.
 \item \emph{CFG-duality ghost.} The CFG 2026 chain-level
 $E_3$-observables of topological Chern--Simons transversely shadow as
 $E_2$-structure on boundary observables (Kapustin--Witten).  Our
 chiral avatar shows the same shadow arises from Dunn--Lurie additivity
 after the Stage-2 specialisation
 $\mathrm{Sp}^{\mathrm{ch}}_{\Sigma_2, C}$; KW corresponds to the
 choice $(\Sigma_2, C) = (\text{topological 2-cycle},
 \text{holomorphic line})$.
\end{enumerate}

Each of (G1)--(G5) is proved as a theorem in the chain-level lane
below and, where it is genuinely stated only at the $(\infty, 1)$-level,
labelled as such.

\subsection*{Stable core}

Fix a saturated CY$_3$ category $\cC$ satisfying H1--H4 of
Theorem~\ref{thm:cy-to-chiral-d3} and an admissible specialisation
$(\Sigma_2, C)$.  Write
\[
 \cF_\cC \;\coloneqq\; \PhiFA_3(\cC) \;\in\; E_3\text{-}\HolFA(X),
 \qquad
 A \;\coloneqq\; \Phi_3^{(\Sigma_2, C)}(\cC)
 \;=\; \mathrm{Sp}^{\mathrm{ch}}_{\Sigma_2, C}(\cF_\cC)
 \;\in\; E_1\text{-}\ChirAlg(C).
\]
The stable core is the following five-line package.  Each line is a
chain-level statement with explicit dg-model, not a slogan.

\begin{center}
\renewcommand{\arraystretch}{1.3}
\begin{tabular}{rll}
 \toprule
 & \textbf{Line} & \textbf{Status} \\
 \midrule
 (C1) & Dunn--Lurie decomposition of $\cF_\cC$: there exists a
   contractible space of factorisations
 & chain-level (Lurie HA 5.1.2.2) \\
   & $\cF_\cC \simeq \cF^{E_1}_\cC \otimes_{E_0} \cF^{E_2}_\cC$ &
   after Stage-1 framing. \\
 (C2) & $\cF^{E_1}_\cC \mapsto A$ under
   $\mathrm{Sp}^{\mathrm{ch}}_{\Sigma_2, C}$; & chain-level \\
   & $\cF^{E_2}_\cC \mapsto Z^{\mathrm{der}}(\Rep^{E_1}(A))$
   under the same. & (Ayala--Francis fact.\ homology). \\
 (C3) & $Z^{\mathrm{der}}(\Rep^{E_1}(A))
   = \RHom_{A \otimes A^{\mathrm{op}}}(A, A)$ in the
   factorisation-category & chain-level \\
   & dg-model; canonical $E_2$-structure by Deligne--Kontsevich. &
   (Lurie HA 5.3.1.14). \\
 (C4) & The $E_2$-braiding admits the explicit chain-level presentation
 & chain-level (Maulik--Okounkov + \\
   & $R^{\mathrm{ch}}(u) \;=\; \Res_{u = u_\star}
   \phi^+_{\mathrm{UV}}(u)$, MO axiom $=$ cocycle condition on
   $\phi^+_{\mathrm{UV}}$. & CoHA shuffle). \\
 (C5) & In the three verified cases, the transverse $E_2$-shadow
   identifies with the &
   chain-level for $\C^3$, K3-fibred, \\
   & representation category of a quasi-triangular Hopf algebra
   $H_\cC$ with explicit $R$-matrix. & toric CY$_3$ (conifold, LP$^2$).\\
 \bottomrule
\end{tabular}
\end{center}

We now prove each line, with the five attack cycles documenting the
residual obstructions and the healing step.


%%% ---------------------------------------------------------------
\subsection*{Attack 1.  Operadic decomposition is not unique}
%%% ---------------------------------------------------------------

\paragraph{Attack.}  Dunn--Lurie gives $E_3 \simeq E_1 \otimes_{E_0}
E_2$, but the decomposition is not canonical: it depends on a choice of
$2$-framing on the little-$3$-discs operad, and different framings
produce $E_1 \otimes E_2$-factor pairs that are quasi-isomorphic but
not equal.  For the hFA $\cF_\cC$ on the CY$_3$ variety, the
$2$-framing datum is the choice of a transverse $2$-plane field; that
is extra data not supplied by the CY$_3$ input alone.  Hence the
supposed ``canonical'' transverse $E_2$-shadow is not canonical, and
any $R$-matrix it produces is unique only up to a $\mathrm{GRT}_1(\Q)$
torsor.

\paragraph{Ghost theorem.}  The space of $E_1 \otimes E_2$ Dunn--Lurie
decompositions of an $E_3$-factorisation algebra is a torsor under the
Grothendieck--Teichm\"uller group $\mathrm{GRT}_1(\Q)$; the
$E_2$-shadow is well-defined on the \emph{quotient} $\cC \sslash
\mathrm{GRT}_1$ but not on $\cC$ itself.

\paragraph{Heal.}  The objection is correct in isolation but
over-counts the choice.  Two observations.

\emph{First}: at the level of representation categories (not
$E_2$-algebras themselves), the Dunn--Lurie decomposition is
\emph{canonical up to contractible choice} -- the braided monoidal
structure on $\Rep^{E_2}(Z^{\mathrm{der}}(\Rep^{E_1}(A)))$ is unique,
even though its $E_2$-algebra model carries a $\mathrm{GRT}_1$-torsor
of presentations.  This is the content of Lurie HA \S5.3.1.14
(Deligne conjecture for dg-categories): the braided monoidal category
$\cZ(\Rep^{E_1}(A))$ is canonical; the specific $E_2$-cochain model
$\RHom_{A \otimes A^{\mathrm{op}}}(A, A)$ is canonical up to
$\mathrm{GRT}_1$.

\emph{Second}: on the CY$_3$ side, the Costello--Li propagator
\cite[\S 6]{CostelloLi2016} fixes a specific associator in the
$\mathrm{GRT}_1$-torsor -- the Kontsevich associator.  Hence the
$R$-matrix is pinned up to contractible choice by the CY$_3$ geometric
input, agreeing with the KT-formality-and-propagator discussion of
Agent~2's chapter.  The canonical object is the braided monoidal
representation category; the $\mathrm{GRT}_1$ ambiguity is a moduli
of presentations and does not affect the Yang--Baxter content.

\paragraph{Surviving core (C1).}  For any Dunn--Lurie decomposition
$\cF_\cC \simeq \cF^{E_1}_\cC \otimes_{E_0} \cF^{E_2}_\cC$ with
associator fixed by the Costello--Li propagator, the pair is
canonically and functorially associated to $\cC$; the resulting
transverse $E_2$-shadow depends only on $\cC$ (and the Costello--Li
framing).


%%% ---------------------------------------------------------------
\subsection*{Attack 2.  Drinfeld centre vs derived centre}
%%% ---------------------------------------------------------------

\paragraph{Attack.}  ``The $E_2$-factor lives on $Z^{\mathrm{der}}
(\Rep^{E_1}(A))$'' conflates the Drinfeld centre (monoidal centre via
half-braidings) with the derived centre (Hochschild cochains); these
are AP-CY4-distinct objects.  Which is it?

\paragraph{Ghost theorem.}  There are two different $E_2$-algebras
naturally associated to the $E_1$-algebra $A$: the derived centre
$Z^{\mathrm{der}}(A) = \RHom_{A \otimes A^{\mathrm{op}}}(A, A)$, which
is an $E_2$-algebra by Deligne--Kontsevich; and the Drinfeld centre
$\cZ(\Rep^{E_1}(A))$, which is a braided monoidal category (hence
$E_2$-\emph{category}) by Joyal--Street--Majid.

\paragraph{Heal.}  The two objects coincide \emph{after passing to
representation categories}.  Specifically, Ben-Zvi--Francis--Nadler
\cite[Thm.~4.1]{BZFN2010} establish, for a locally presentable stable
$\infty$-category $\cC$ dualisable over $\cC \otimes \cC^{\mathrm{op}}$,
the identification
\[
 \mathrm{HH}^\bullet_{\mathrm{cat}}(\cC, \cC) \;=\;
 \End_{\cC \otimes \cC^{\mathrm{op}}}(\cC) \;\simeq\; Z(\cC)
\]
as $E_2$-algebras in $\mathrm{Pr}^{\mathrm{L}}$; the first is the
categorical Hochschild cochains and the third is the Drinfeld centre.
Specialising to $\cC = \Rep^{E_1}(A)$ with $A$ a chiral algebra along
$C$ and dualisable over $A \otimes A^{\mathrm{op}}$, we obtain
\[
 \underbrace{Z^{\mathrm{der}}(\Rep^{E_1}(A))}_{\text{derived Drinfeld
 centre on the chiral side}}
 \;\simeq\;
 \underbrace{\RHom_{A \otimes A^{\mathrm{op}}}(A, A)}_{\text{Hochschild
 cochains of $A$ as a dg-algebra object}}
 \;\simeq\;
 \underbrace{\cZ(\Rep^{E_1}(A))}_{\text{Drinfeld centre of the
 representation category}},
\]
all three as $E_2$-algebras.  The distinction between the two ghost
candidates is a notational artifact: they are the same $E_2$-object
when the representation category is sufficiently dualisable.  What we
call $Z^{\mathrm{der}}(\Rep^{E_1}(A))$ throughout is this common
$E_2$-algebra.

\paragraph{Spectral-parameter refinement (BZFN + Vol~III cache
entries 46--50).}  One subtlety: for chiral algebras along $C$, the
chiral-Hochschild complex $\mathrm{ChirHoch}(A, A)$ and the
categorical Hochschild complex $\mathrm{HH}^\bullet_{\mathrm{cat}}
(\Rep^{E_1}(A))$ differ by the evaluation-module data.  The spectral
parameter $u$ entering the $R$-matrix comes from the evaluation modules
$\mathrm{ev}_u \colon A \to \End(V_u)$, not from the centre itself.
The derived centre has no spectral parameter; the spectral parameter
enters when one evaluates the universal $R$-matrix
$R^{\mathrm{ch}} \in Z^{\mathrm{der}} \otimes Z^{\mathrm{der}}$ on a
pair of evaluation modules $(V_{u_1}, V_{u_2})$, producing
$R^{\mathrm{ch}}_{V_{u_1}, V_{u_2}}(u_1 - u_2)$.  This aligns exactly
with the spectral-parameter discussion in
\texttt{appendices/first\_principles\_cache.md}: spectral parameter
from evaluation modules, not from the centre.

\paragraph{Surviving core (C2)--(C3).}  The second factor
$\cF^{E_2}_\cC$ of the Dunn--Lurie decomposition, after the Stage-2
specialisation, yields the $E_2$-algebra $Z^{\mathrm{der}}(\Rep^{E_1}
(A))$, and this $E_2$-algebra coincides with the Hochschild cochain
complex of $A$ viewed as a dg-algebra object in
$E_1$-$\ChirAlg(C)$.  The representation category of this $E_2$-algebra
carries a canonical braided monoidal structure (the Drinfeld-centre
braiding), and on evaluation modules the braiding evaluates to a
holomorphic $R$-matrix $R^{\mathrm{ch}}(u)$.


%%% ---------------------------------------------------------------
\subsection*{Attack 3.  The R-matrix as a residue: but of what?}
%%% ---------------------------------------------------------------

\paragraph{Attack.}  The Maulik--Okounkov $R$-matrix is constructed
from the K\"unneth theorem on stable envelopes on Nakajima varieties.
Claiming it is a residue of a UV gluing cocycle $\phi^+_{\mathrm{UV}}$
is either wrong (MO is not a residue, it is a transition matrix) or
vacuous (every transition matrix across a wall is a residue of
something by Cauchy's theorem).

\paragraph{Ghost theorem.}  For a CoHA $\cH$ associated to a
Nakajima quiver variety $X^+$ equivariant under a torus $T$, there
exists a meromorphic $\End(\cH \otimes \cH)$-valued $1$-form
$\phi^+_{\mathrm{UV}}(u)$ on the equivariant moduli
$\Spec H^*_T(\mathrm{pt})$, regular on chambers and with simple poles
along the chamber walls, such that
\[
 R^{\mathrm{MO}}(u_1 - u_2)
 \;=\;
 \mathrm{Res}_{u_1 - u_2 = u_\star}
 \phi^+_{\mathrm{UV}}(u_1 - u_2),
\]
and the Yang--Baxter relation plus unitarity are the
$\mathrm{Cech}$-cocycle condition for $\phi^+_{\mathrm{UV}}$ on the
triple-wall codimension-two intersection.

\paragraph{Heal.}  We construct $\phi^+_{\mathrm{UV}}$ explicitly.  Let
$X^+ = \cM^+_{\mathrm{eff}}(\cC)$ denote the effective moduli of
positive-stability objects of $\cC$ (toric quiver + potential for
toric CY$_3$; derived moduli of sheaves for K3 $\times$ E).  The
CoHA is
\[
 Y^+(\cC) \;=\; H^\bullet_T(X^+, \phi_W)
\]
with Hall product from the stack of short exact sequences (Kontsevich--
Soibelman \cite{KontsevichSoibelmanCoHA}; Schiffmann--Vasserot
\cite{SchiffmannVasserot2013} for the quiver case).  Write $\chi$ for
the torus character lattice and $\mathfrak{t}^*_\C =
\chi \otimes_\Z \C$ for its complexification.  A chamber of the
equivariant moduli is a connected component of the complement of the
wall arrangement $\cW \subset \mathfrak{t}^*_\C$, where the walls are
the loci $\{u \in \mathfrak{t}^*_\C : \alpha(u) = 0\}$ for $\alpha$
running over the $T$-characters of the normal bundles to fixed-point
components of $X^+$.

The \emph{UV positive-half gluing cocycle} is the meromorphic
$1$-form on $\mathfrak{t}^*_\C$ with values in $\End(Y^+ \otimes Y^+)$
defined on a chamber $\mathcal{K}$ by
\[
 \phi^+_{\mathrm{UV}}(u) \Big|_{\mathcal{K}}
 \;=\;
 \mathrm{Stab}_{\mathcal{K}}^+
 \circ (\mathrm{Stab}_{\mathcal{K}}^+)^{-1}
 \;+\;
 \sum_{\text{walls } \alpha} \frac{r_\alpha \, du}{\alpha(u)},
\]
where $\mathrm{Stab}_{\mathcal{K}}^+$ is the Maulik--Okounkov stable
envelope associated to the chamber $\mathcal{K}$ (taking values in
the fixed-point sector), and $r_\alpha$ is the residue at the wall
$\alpha$ computed from the hyperbolic-restriction functor across that
wall.  Concretely: for two adjacent chambers $\mathcal{K}_1,
\mathcal{K}_2$ separated by a single wall $\alpha = 0$, the
Maulik--Okounkov $R$-matrix $R^{\mathrm{MO}}(\alpha)$ across the wall
is the transition matrix between the two stable envelopes,
\[
 \mathrm{Stab}_{\mathcal{K}_2}^+ \;=\;
 R^{\mathrm{MO}}(\alpha) \cdot \mathrm{Stab}_{\mathcal{K}_1}^+,
\]
which is precisely the residue of $\phi^+_{\mathrm{UV}}$ at that wall
in the direction from $\mathcal{K}_1$ to $\mathcal{K}_2$.

\paragraph{Cocycle condition = Yang--Baxter.}  Across a codimension-two
wall, three chambers meet: $\mathcal{K}_1, \mathcal{K}_2, \mathcal{K}_3$
with pairwise transitions $R^{\mathrm{MO}}_{12}, R^{\mathrm{MO}}_{23},
R^{\mathrm{MO}}_{13}$.  The cocycle condition reads
$R^{\mathrm{MO}}_{13} = R^{\mathrm{MO}}_{23} \cdot R^{\mathrm{MO}}_{12}$;
for three interacting particles with charges summing to zero on a
single codimension-two wall, the condition
$R^{\mathrm{MO}}_{13}(u_1 + u_2) = R^{\mathrm{MO}}_{23}(u_2) \cdot
R^{\mathrm{MO}}_{12}(u_1)$ is the Yang--Baxter equation in the
reducible form; using the analytic structure (meromorphic
continuation of $R^{\mathrm{MO}}$ is rational in $u$), one passes to
the braided form
\[
 R^{\mathrm{MO}}_{12}(u_1 - u_2) \, R^{\mathrm{MO}}_{13}(u_1 - u_3) \,
 R^{\mathrm{MO}}_{23}(u_2 - u_3) \;=\;
 R^{\mathrm{MO}}_{23}(u_2 - u_3) \, R^{\mathrm{MO}}_{13}(u_1 - u_3) \,
 R^{\mathrm{MO}}_{12}(u_1 - u_2).
\]
This is Maulik--Okounkov \cite[Theorems~4.6.1 and 5.1.4]{MaulikOkounkov2012}
re-read as a $\mathrm{Cech}$-cocycle condition on
$\phi^+_{\mathrm{UV}}$.

\paragraph{Unitarity = cocycle-reversal.}  Exchanging two adjacent
chambers reverses the order of the particles; the resulting
$R^{\mathrm{MO}}_{21}(-u) \cdot R^{\mathrm{MO}}_{12}(u) = \mathrm{id}$
is the cocycle reversal $\phi^+_{\mathrm{UV}} + \bar\phi^+_{\mathrm{UV}}
= d(\cdots)$, where $\bar\phi^+_{\mathrm{UV}}$ is the chamber-reversed
cocycle.

\paragraph{Surviving core (C4).}  The $R$-matrix of the transverse
$E_2$-shadow is
\[
 R^{\mathrm{ch}}(u) \;=\; R^{\mathrm{MO}}(u)
 \;=\;
 \Res_{v = u_\star} \phi^+_{\mathrm{UV}}(v),
 \qquad
 v = u_1 - u_2,
\]
with $\phi^+_{\mathrm{UV}}$ the UV positive-half gluing cocycle
constructed from the stable-envelope data.  The Yang--Baxter equation
and unitarity are the chain-level cocycle and cocycle-reversal
conditions on $\phi^+_{\mathrm{UV}}$.  The hexagon axioms of the
braided monoidal structure on
$\Rep^{E_2}(Z^{\mathrm{der}}(\Rep^{E_1}(A)))$ are the corresponding
codimension-three coherences.


%%% ---------------------------------------------------------------
\subsection*{Attack 4.  Hexagon and the Drinfeld-double presentation}
%%% ---------------------------------------------------------------

\paragraph{Attack.}  Chain-level existence of an $R$-matrix with YBE
is weaker than the statement that $Z^{\mathrm{der}}(\Rep^{E_1}(A))$
is the representation category of a quasi-triangular Hopf algebra.
For Hopf-algebraic input one needs: coproduct $\Delta$, antipode $S$,
the six hexagon axioms, and the identification of the $R$-matrix with
the action of a universal element in $H \otimes H$.  Without these,
the braided structure is a generic $E_2$-algebra with no Hopf
presentation.

\paragraph{Ghost theorem.}  For any CY$_3$ category $\cC$ admitting a
CoHA $Y^+(\cC)$ with locally-finite graded pieces, the transverse
$E_2$-shadow of $\cC$ is the representation category of the
Drinfeld double $H_\cC \coloneqq D(Y^+(\cC)) = Y^+(\cC) \bowtie
Y^-(\cC)$, a quasi-triangular Hopf algebra with universal $R$-matrix.

\paragraph{Heal.}  For the three verified CY$_3$ cases we construct the
Drinfeld double presentation explicitly.

\emph{Case I: $\cC = \Perf(\C^3)$ (toric CY$_3$, no compact 4-cycles).}
The CoHA is $Y^+(\widehat{\mathfrak{gl}}_1)$, the positive half of the
affine Yangian of $\mathfrak{gl}_1$ (Schiffmann--Vasserot 2013).  Its
graded dimension is the MacMahon series $M(q) = \prod_{n \geq 1}
(1-q^n)^{-n}$.  The Drinfeld double
$D(Y^+(\widehat{\mathfrak{gl}}_1))$ is the full affine Yangian
$Y(\widehat{\mathfrak{gl}}_1)$, by Schiffmann--Vasserot
\cite{SchiffmannVasserot2013} Theorem~8.3 (positive and negative halves
are graded duals with reversed coproduct, and the Drinfeld double of
a positive half recovers the Yangian).  The universal $R$-matrix is
Drinfeld's New Realisation $R \in Y \otimes Y$ with $R = 1 +
\hbar \sum_\alpha E_\alpha \otimes F_\alpha + O(\hbar^2)$ in the
PBW basis; its evaluation on Fock spaces of charge $(n_1, n_2)$ is
\[
 R^{\mathrm{ch}}_{\mathrm{Fock}_{n_1}, \mathrm{Fock}_{n_2}}(u) \;=\;
 \prod_{\square_1 \in \lambda_1} \prod_{\square_2 \in \lambda_2}
 \frac{u - u_\star(\square_1, \square_2)}{u + u_\star(\square_1,
 \square_2)}
\]
with $u_\star(\square_1, \square_2) = h_1 (c_1 - c_2) +
h_2 (r_1 - r_2)$ the two-Omega-background wall-crossing point
(Feigin--Jimbo--Miwa--Mukhin \cite{FJMM2017}).  This is the
matrix-Fock-module $R$-matrix of the affine Yangian.  The hexagon
axioms are the bialgebra compatibility for the positive-half coproduct
from deconcatenation on $Y^+(\widehat{\mathfrak{gl}}_1)$ extended to
the double.

\emph{Case II: $\cC = D^b(\Coh(K3 \times E))$ (compact CY$_3$, reduced
K3-fibred branch).}  The reduced Mukai-lattice structure forces the
transverse $E_2$-shadow into the super-Yangian family: the pairing
$(\widetilde H^*_+(K3), \chi_{\mathrm{Muk}})$ has signature $(4, 20)$
with the $(1,1)$-part $\widetilde H^{1,1}(K3) = 20$-dim giving the
bosonic sector and the $(1,0) + (0,1) + (2,0) + (0,2)$-part $= 4$-dim
giving the fermionic sector (Mukai's Hodge involution on
$\widetilde H^*$).  The CoHA carries the $\mathbb Z/2$-graded
super-structure inherited from the $+1/-1$-eigenspaces of this Hodge
involution, and the Drinfeld double is
\[
 H_{K3 \times E} \;=\; D(Y^+(\cC_{K3 \times E})) \;=\;
 Y_{\mathfrak{osp}}(4|20)
 \;=\; Y(\mathfrak{gl}(4|20))^\theta,
\]
the fixed subalgebra of $Y(\mathfrak{gl}(4|20))$ under the Chevalley
involution $\theta$ preserving the Mukai-orthogonal symplectic form.
The signature $(4, 20)$ is forced rigidly by the Hodge decomposition
of $\widetilde H^*(K3, \C)$; the super-Yangian assignment agrees with
the Vol~III cache memory entry \emph{``K3 super-Yangian
$Y_{\mathfrak{osp}}(4|20)$''}
(\texttt{project\_super\_yangian\_Y\_osp\_4\_20.md}).

\emph{Case III: toric CY$_3$ with a compact 4-cycle (LP$^2$,
conifold).}  The CoHA is given by Rapcak--Soibelman--Yang--Zhao
\cite{RSYZ} as a shuffle algebra; the Drinfeld double is the
quantum toroidal algebra $U_{q, t}(\widehat{\widehat{\mathfrak{gl}}}_r)$
with $r$ depending on the fan combinatorics.  For LP$^2$, $r = 1$ and
the double is the quantum toroidal $\mathfrak{gl}_1$ at a specific
parameter specialisation.  For the conifold, $r = 2$ and the double
is a Drinfeld-deformed symmetric-pair toroidal algebra.

\paragraph{Hexagon verification.}  In each case, the hexagon axioms
are the Drinfeld double's structure theorem
(Majid \cite{Majid1990}, Kassel \cite{KasselQuantumGroups} Ch.~IX,
Etingof--Kazhdan \cite{EtingofKazhdan2008} Part~IV for the infinite-
dimensional case).  On the chain level, the hexagon identities
\[
 R_{12, 3} = R_{13} \cdot R_{23}, \qquad R_{1, 23} = R_{13} \cdot
 R_{12}
\]
restricted to representations of the Drinfeld double are tautologies,
being part of the definition of quasi-triangularity; on the
$E_2$-algebra level they are the iterated-insertion coherences of the
Fulton--MacPherson operad $\FM_n(\C)$ acting on
$B_{E_2}(Z^{\mathrm{der}})$.  The passage between the two
presentations (Hopf-algebraic hexagons $\Leftrightarrow$
$E_2$-operadic coherences) is Lurie HA \S5.3.1.14.

\paragraph{Surviving core (C5).}  For each of the three verified
cases, the transverse $E_2$-shadow is the representation category of a
quasi-triangular Hopf algebra $H_\cC$ whose $R$-matrix is the MO
$R$-matrix, equivalently the residue of $\phi^+_{\mathrm{UV}}$.  For
more general compact CY$_3$ (quintic, Calabi--Yau 3-folds admitting
only a derived-moduli construction with no known CoHA) the transverse
$E_2$-shadow exists as an abstract $E_2$-algebra but the Hopf
presentation is the CY-B$_3$ conjectural content.


%%% ---------------------------------------------------------------
\subsection*{Attack 5.  CFG 2026 correspondence and the Kapustin--Witten face}
%%% ---------------------------------------------------------------

\paragraph{Attack.}  The CFG 2026 theorem states that the $E_3$-observables
of topological Chern--Simons on $\R^3$ are Koszul-dual to
$U_q(\mathfrak g)$-mod, an $E_2$-category.  Applying the same logic to
hCS on $\C^3$ yields an $E_2$-shadow on boundary observables
(Kapustin--Witten).  But the chiral setting has holomorphic (not
topological) data; the boundary observables are $E_1$-chiral on a
curve $C \subset \C^3$.  Why should our transverse $E_2$-shadow
coincide with CFG's $E_2$-shadow?

\paragraph{Ghost theorem.}  The transverse $E_2$-shadow of
$\PhiFA_3(\cC)$ after the Stage-2 specialisation equals the
CFG $E_2$-shadow of topological Chern--Simons after choosing
$(\Sigma_2, C) = (\text{topological 2-cycle}, \text{holomorphic line})$;
the two constructions are the same $E_2$-algebra under the Langlands
duality exchange of topological and holomorphic data.

\paragraph{Heal.}  The correspondence is not an accident.  Both are
instances of the same Dunn--Lurie decomposition, applied with different
input framings.

\emph{CFG side.}  CFG 2026 studies $E_3$-observables of topological
Chern--Simons on a framed $3$-manifold $M^3$.  The Dunn--Lurie
decomposition $E_3 \simeq E_1 \otimes E_2$ with $(\Sigma_2, C) =
(\text{topological plane}, \text{topological line})$ consumes the
$E_1$-factor by evaluating along $C$ (Wilson-line observables) and
leaves the $E_2$-factor on transverse directions; the resulting
$E_2$-algebra is $U_q(\mathfrak g)$-mod.

\emph{Our side.}  We study $E_3$-hFA $\PhiFA_3(\cC)$ on a CY$_3$
variety $X$.  The Dunn--Lurie decomposition with
$(\Sigma_2, C) = (\text{transverse holomorphic } 2\text{-cycle},
\text{holomorphic curve})$ consumes the $E_1$-factor by Stage-2
specialisation along $C$ (chiral factorisation along the curve) and
leaves the $E_2$-factor on the transverse $2$-cycle; the resulting
$E_2$-algebra is $Z^{\mathrm{der}}(\Rep^{E_1}(A))$, which is
$U_q(\mathfrak g)$-mod-like for Kac--Moody-type CoHA (Schiffmann--
Vasserot, $\mathfrak g = \widehat{\mathfrak{gl}}_1$).

\emph{Identification.}  Under the Costello--Gaiotto--Yagi 5D hCS-to-
Yangian-VOA theorem (simply-laced $\mathfrak g$, convergent to all
orders), the hCS gauge theory on $\R \times \C^2$ with Wilson-line
observables along $C = \R$ is a holomorphic deformation of
topological Chern--Simons on $\R^3$: the chiral $E_1$-direction along
$C$ replaces the topological one, and the transverse $\C^2 \simeq \R^4$
is the holomorphic analogue of the topological $\R^2$.  The
Dunn--Lurie decompositions of the two ambient $E_3$-objects are
intertwined by this gauge-theoretic identification, and the two
transverse $E_2$-shadows agree.  Concretely: CFG's $U_q(\mathfrak g)$-mod
is the \emph{rational} evaluation on spectral-parameter infinity of
our $R^{\mathrm{ch}}(u)$; our $R^{\mathrm{ch}}(u)$ is the
holomorphic deformation of CFG's $R^{\mathrm{CFG}}$ away from infinity.

\paragraph{Surviving core (G5).}  The CFG 2026 $E_2$-shadow is the
specialisation of our transverse $E_2$-shadow at
$(\Sigma_2, C) = (\text{topological plane}, \text{topological line})$;
our transverse $E_2$-shadow is the holomorphic (CY$_3$) refinement of
CFG at generic $(\Sigma_2, C) = (\text{holomorphic } 2\text{-cycle},
\text{holomorphic curve})$.  Both are instances of Dunn--Lurie applied
to an $E_3$-hFA; the difference is in the framing datum, not in the
categorical apparatus.


%%% ---------------------------------------------------------------
\subsection*{Re-attack 1.  The decomposition canonicity, revisited}
%%% ---------------------------------------------------------------

\paragraph{Re-attack.}  The Kontsevich--Tamarkin formality theorem
pins the Dunn--Lurie decomposition of the \emph{categorical}
$E_3$-structure on $\HH^\bullet(\cC, \cC)$ up to
$\mathrm{GRT}_1(\Q)$; the Costello--Li propagator selects the
Kontsevich-associator point.  But the CY$_3$ trace class
$\eta \in \HH_{-3}(\cC)$ is extra structure, and it should impose
further compatibility on the Dunn--Lurie decomposition.  Does the
trace pin the decomposition beyond the $\mathrm{GRT}_1$ torsor, or is
the $\mathrm{GRT}_1$ ambiguity genuine?

\paragraph{Re-heal.}  The CY$_3$ trace imposes \emph{cyclic
compatibility} on the Dunn--Lurie decomposition: the $E_1$-factor
must be cyclic $A_\infty$ with trace $\eta$, and the $E_2$-factor
must carry a cyclic structure of degree $3 - 1 = 2$ (shifted by the
CY-shift, Kontsevich--Soibelman \cite{KontsevichSoibelmanDefI}).  This
cyclic compatibility cuts down the $\mathrm{GRT}_1$ torsor of
decompositions to its \emph{cyclic} subgroup $\mathrm{GRT}_1^{\mathrm{cyc}}
(\Q) \subset \mathrm{GRT}_1(\Q)$; the Kontsevich associator is in the
cyclic subgroup (Tamarkin \cite{TamarkinDGCat}), so the
Costello--Li selection is consistent.  For non-cyclic algebras (not
supplied by a CY category), the ambiguity would be the full
$\mathrm{GRT}_1$.

Residual frontier: the cyclic $\mathrm{GRT}_1$ is not trivial; its
$\pi_0$-class is the obstruction to strict $2$-shifted cyclicity on
the $E_2$-factor.  At chain level this is the Costello--Gwilliam--Li
$\hbar$-loop anomaly on the $E_2$-factor; at $d = 3$ it vanishes for
$\mathfrak g$ simply-laced on the $\C^3$-specialisation, giving the
all-orders theorem.


%%% ---------------------------------------------------------------
\subsection*{Re-attack 2.  E_2 on Rep vs E_2 on Z^der}
%%% ---------------------------------------------------------------

\paragraph{Re-attack.}  We identified
$Z^{\mathrm{der}}(\Rep^{E_1}(A)) \simeq \cZ(\Rep^{E_1}(A))$ as
$E_2$-algebras via Ben-Zvi--Francis--Nadler; that gives the
$E_2$-structure on the \emph{centre}.  But the target of
\emph{representations} is an $E_2$-category, not an $E_2$-algebra.
Have we conflated these?

\paragraph{Re-heal.}  The $E_2$-category of representations of an
$E_2$-algebra $Z$ is $\Rep^{E_2}(Z)$; by Lurie HA \S5.5.1 this is
a braided monoidal stable $\infty$-category, and the braided
monoidal structure is by definition the Drinfeld-centre structure
induced from $\Rep^{E_1}(A)$ via passing to the centre.  There is no
conflation: the $E_2$-algebra
$Z^{\mathrm{der}} = \RHom_{A \otimes A^{\mathrm{op}}}(A, A)$ is
an algebraic object; its $E_2$-category $\Rep^{E_2}(Z^{\mathrm{der}})$
is a categorical object; and the identification
$\Rep^{E_2}(Z^{\mathrm{der}}(\Rep^{E_1}(A))) \simeq
\cZ(\Rep^{E_1}(A))$ is Lurie HA \S5.3.1.14 combined with
Ben-Zvi--Francis--Nadler.

The four-way identification is then:
\[
 \underbrace{\RHom_{A \otimes A^{\mathrm{op}}}(A, A)}_{E_2\text{-algebra,
 chain-level}}
 \;\leadsto\;
 \underbrace{\Rep^{E_2}\!\bigl(\RHom_{A \otimes A^{\mathrm{op}}}(A, A)
 \bigr)}_{E_2\text{-category}}
 \;\simeq\;
 \underbrace{\cZ(\Rep^{E_1}(A))}_{\text{Drinfeld centre, categorical}}
 \;\simeq\;
 \underbrace{\Rep(D(H_\cC))}_{\text{Hopf-algebraic}},
\]
the last with $H_\cC$ the quasi-triangular Hopf algebra of the
transverse $E_2$-shadow (when it exists; for the three verified cases).


%%% ---------------------------------------------------------------
\subsection*{Re-attack 3.  The chamber structure depends on the CY_3}
%%% ---------------------------------------------------------------

\paragraph{Re-attack.}  The UV gluing cocycle $\phi^+_{\mathrm{UV}}$
was defined via the equivariant chamber structure on
$\mathfrak{t}^*_\C$.  For a general CY$_3$ without a torus action, the
chamber structure is absent; our residue presentation then has no
direct analogue.  Is the transverse $E_2$-shadow still well-defined?

\paragraph{Re-heal.}  Yes.  For CY$_3$ without equivariance, the MO
$R$-matrix is replaced by Okounkov's K-theoretic stable envelope
on derived moduli (Okounkov \cite{Okounkov2017Lectures} \S 8--\S 9),
valued in the K-theory of the pair $(\cM^+, \cM^+_{\mathrm{fixed}})$.
The residue presentation then reads
\[
 R^{\mathrm{ch}}(u) \;=\;
 \mathrm{Res}_{\mathrm{wall}} \phi^+_{\mathrm{UV,K}}(u),
\]
with $\phi^+_{\mathrm{UV,K}}$ a K-theoretic $1$-form on the
$\hbar$-moduli.  For $\hbar \to 0$ this reduces to the cohomological
MO construction; for $\hbar$ generic it produces the elliptic or
trigonometric $R$-matrix (elliptic, Okounkov--Rimanyi).

Four equivariance strata (Vol~III cache):
(i)  toric $T^d$: cohomological MO, residue at toric fixed point
(ii) reduced $\C^\times + \mathrm{Aut}(X)$: parameter-free MO
(iii) orbifold inertia $I(X/G)$: McKay equivariant MO
(iv) lattice-polarised period domain: Borcherds lift
The residue presentation adapts to each stratum; only the
$\phi^+_{\mathrm{UV}}$ changes analytic form.


%%% ---------------------------------------------------------------
\subsection*{Re-attack 4.  K3 x E signatures: (4,20) vs product factors}
%%% ---------------------------------------------------------------

\paragraph{Re-attack.}  For $K3 \times E$, we asserted
$H_{K3 \times E} = Y_{\mathfrak{osp}}(4|20)$.  But $K3 \times E$ is a
product of a CY$_2$ factor (K3) and a $d=1$ factor (elliptic $E$); by
Dunn--Lurie on the product, the CoHA should factor as
$Y^+(K3) \otimes Y^+(E)$ with $Y^+(E)$ a lattice VOA.  This factored
form is not manifestly $Y_{\mathfrak{osp}}(4|20)$.

\paragraph{Re-heal.}  The decomposition is at the level of
$\mathrm{CoHA}$, not at the level of the Drinfeld double $H_{K3 \times E}$.
On K3 alone (CY$_2$), the CoHA is
$Y^+(\widetilde H^*(K3)) \otimes_k k$, and the Drinfeld double of
this CY$_2$ CoHA is the Mukai-braided enhancement, not the super-Yangian.
The super-structure and the $(4, 20)$-signature appear after the
Mukai-Hodge involution is read as an $\mathrm{osp}$-symmetry breaking;
this is a feature of the $d = 3$ total space (via the elliptic fibre
as spectral-parameter generator), not of the K3 CY$_2$ factor alone.

Concretely: the elliptic fibre $E$ supplies the spectral parameter $u$
(holomorphic along $E$); the K3 base supplies the $R$-matrix target
space $\widetilde H^*(K3) = \C^{4|20}$.  The pair
$(\widetilde H^*(K3), E\text{-spectral parameter})$ is an evaluation
module for the Yangian $Y(\mathfrak{gl}(4|20))$; the Mukai involution
cuts this down to $Y(\mathfrak{gl}(4|20))^\theta = Y_{\mathfrak{osp}}
(4|20)$.  The product decomposition of the CoHA at $d = 2$ becomes
the evaluation-representation data of the super-Yangian at $d = 3$;
the two statements are compatible, living at different layers of the
Dunn--Lurie decomposition.


%%% ---------------------------------------------------------------
\subsection*{Re-attack 5.  Residual conjectural content}
%%% ---------------------------------------------------------------

\paragraph{Re-attack.}  After the five healing steps, what remains
conjectural?

\paragraph{Status table.}
\begin{center}
\renewcommand{\arraystretch}{1.3}
\begin{tabular}{lll}
 \toprule
 \textbf{Statement} & \textbf{Status} & \textbf{Residue} \\
 \midrule
 (C1) Dunn--Lurie decomposition up to $\mathrm{GRT}_1^{\mathrm{cyc}}$ &
 chain-level theorem &
 Lurie HA 5.1.2.2 + Tamarkin \\
 (C2) Stage-2 sends $\cF^{E_2}$ to $Z^{\mathrm{der}}(\Rep^{E_1}(A))$ &
 chain-level theorem &
 Ayala--Francis \S 3.1 \\
 (C3) $Z^{\mathrm{der}} = \RHom_{A\otimes A^{\mathrm{op}}}(A, A)$ as
 $E_2$-algebra &
 chain-level theorem &
 Lurie HA 5.3.1.14 \\
 (C4) $R^{\mathrm{ch}}(u) = \Res \phi^+_{\mathrm{UV}}(u)$, YBE from
 cocycle & chain-level theorem &
 MO 2012 + re-read \\
 (C5) Hopf presentation $H_\cC$ for verified cases &
 chain-level theorem &
 Case-by-case \\
 (C5$^*$) Hopf presentation for compact non-K3-fibred CY$_3$ &
 CY-B$_3$ conjectural &
 quintic, CICY, \ldots \\
 (G5) CFG $E_2$-shadow $=$ Kapustin--Witten specialisation &
 $(\infty, 1)$-theorem &
 Dunn--Lurie compatibility \\
 \bottomrule
\end{tabular}
\end{center}

The genuinely open content is (C5$^*$): the Hopf-algebraic presentation
of the transverse $E_2$-shadow for compact CY$_3$ without a CoHA
construction.  For such CY$_3$, the transverse $E_2$-shadow exists as
an abstract $E_2$-algebra (by C1--C3), and its braiding exists as an
$R$-matrix (by C4, on a suitable derived-moduli stratum); but the
identification with $\Rep$ of a quasi-triangular Hopf algebra is the
CY-B$_3$ conjectural content.  This is compatible with
Remark~\ref{rem:cy-b-d3-precise}(c) and
Conjecture~\ref{conj:kapranov-3shifted-exterior-koszul}.


%%% =====================================================================
\subsection*{Explicit chain-level derivation (detailed)}
%%% =====================================================================

We now compile the explicit chain-level formulas supporting the stable
core lines (C1)--(C5), for the three verified CY$_3$ cases.

\subsubsection*{The Dunn--Lurie tensor and its infinitesimal form}

Lurie HA \S5.1.2.2 states that the little-cubes operad
$E_n$ admits a tensor-decomposition at the level of
$\infty$-operads: for $n = p + q$,
\[
 E_n \;\simeq\; E_p \otimes E_q
\]
as $\infty$-operads; and the $E_p$-algebra on the left-hand side is
precisely an $E_q$-algebra in $E_p$-algebras.  Infinitesimally, the
Chevalley--Eilenberg cochain complex decomposes as
\[
 \mathrm{CE}^\bullet(E_n\text{-algebra})
 \;\simeq\; \mathrm{CE}^\bullet(E_p\text{-algebra in }E_q\text{-algebras}),
\]
with the two differentials $d_{E_p}$ and $d_{E_q}$ commuting (both
are squared-to-zero, and their graded commutator is $d_{E_0}$ = zero
in the stable case).

For our case $n = 3 = 1 + 2$:
\[
 \cF_\cC \;\in\; E_3\text{-}\HolFA(X) \qquad\Longrightarrow\qquad
 \cF_\cC \;\simeq\; \cF^{E_1}_\cC \otimes_{E_0} \cF^{E_2}_\cC
\]
with $\cF^{E_1}_\cC \in E_1\text{-}\HolFA(X)$ (Beilinson--Drinfeld
chiral on the $1$-cycle factor of the framed decomposition of $X$) and
$\cF^{E_2}_\cC \in E_2\text{-}\HolFA(X)$ (fully holomorphic
factorisation on the complementary $2$-cycle).

\subsubsection*{The Stage-2 specialisation}

Stage-2 of the CY-to-chiral functor is the factorisation-homology
specialisation $\mathrm{Sp}^{\mathrm{ch}}_{\Sigma_2, C}$, which,
concretely on the $E_1 \otimes E_2$-factored form, acts by:
\[
 \mathrm{Sp}^{\mathrm{ch}}_{\Sigma_2, C}
 \colon \cF^{E_1}_\cC \otimes_{E_0} \cF^{E_2}_\cC
 \longmapsto
 A \otimes_{E_0} \fZ_\cC,
\]
where $A = \mathrm{Sp}^{\mathrm{ch}}_{\Sigma_2, C}(\cF^{E_1}_\cC)$ is
the chiral algebra along $C$ (the factorisation homology of
$\cF^{E_1}_\cC$ along $\Sigma_2 \setminus C$, pulled back to $C$),
and $\fZ_\cC = \int_{\Sigma_2} \cF^{E_2}_\cC$ is the factorisation
homology of $\cF^{E_2}_\cC$ over the $\Sigma_2$-cycle.

Ayala--Francis \cite[\S 3.1]{AyalaFrancis2015} establish that this
second factor
$\fZ_\cC = \int_{\Sigma_2} \cF^{E_2}_\cC$ identifies as the derived
Drinfeld centre of the representation category of $A$:
\[
 \fZ_\cC \;\simeq\; Z^{\mathrm{der}}(\Rep^{E_1}(A)).
\]
The identification is explicit: both sides compute the derived
endomorphisms of the category $\Rep^{E_1}(A)$ as a module over itself
via the $E_2$-factorisation structure; the Ayala--Francis
factorisation-homology-over-manifolds theorem identifies these on
$(\Sigma_2 = S^2)$-factorisation.

\subsubsection*{The explicit Hochschild cochain model}

For concreteness: in the dg-model, write $A$ as a unital associative
dg-algebra in the symmetric monoidal category $E_1\text{-}\ChirAlg(C)$.
Then the Hochschild cochain complex of $A$ is
\[
 \mathrm{HH}^\bullet(A, A) \;=\;
 \Bigl(\bigoplus_{n \geq 0} \Hom_{E_1}(A^{\otimes n}, A),
 \; b + \mu_* \Bigr),
\]
with $b$ the Hochschild coboundary and $\mu_*$ the chiral-bracket
contribution.  The Deligne--Kontsevich $E_2$-action is via brace
operations $\{-,\{-,\ldots\}\}$ (Gerstenhaber's original brace
\cite{Gerstenhaber}) upgraded to $E_2$-brace (Kontsevich--Soibelman
\cite{KontsevichSoibelmanDefI}; Tamarkin \cite{TamarkinDGCat};
McClure--Smith \cite{McClureSmith}).  The identification
\[
 Z^{\mathrm{der}}(\Rep^{E_1}(A)) \;\simeq\; \mathrm{HH}^\bullet(A, A)
 \qquad \text{as } E_2\text{-algebras}
\]
is Lurie HA \S5.3.1.14 (Deligne conjecture for dg-algebras).

\subsubsection*{The R-matrix as cocycle residue: toric case}

For $\cC = \Perf(\C^3)$, the CoHA is $Y^+(\widehat{\mathfrak{gl}}_1)$
and the UV gluing cocycle is constructed explicitly via torus
equivariance.  Let $T = (\C^\times)^2 \subset (\C^\times)^3$ be the
CY$_3$-preserving subtorus (fixing $h_1 h_2 h_3 = 0$, i.e. acting on
$\C^3$ via $(t_1, t_2, t_3)$ with $t_1 t_2 t_3 = 1$).  The fixed
locus of $T$ on $\Hilb^n(\C^3)$ is the finite set of
torus-fixed ideals $I_\lambda$ indexed by $3d$-partitions $\lambda$
(plane partitions of size $n$).  The stable envelope takes values in
\[
 \mathrm{Stab}_{\mathcal{K}} \;\colon\;
 K_T(\Hilb^n(\C^3)^T) \;\longrightarrow\; K_T(\Hilb^n(\C^3)),
\]
with $\mathrm{Stab}_{\mathcal{K}}([\lambda])$ a class supported on
the attracting set of $I_\lambda$ in the chamber $\mathcal{K}$.
Wall-crossing between two chambers
$\mathcal{K}^\pm$ separated by a wall $\{h_i = 0\}$ gives the
MO $R$-matrix across that wall, with explicit formula (FJMM
\cite{FJMM2017})
\[
 R^{\mathrm{MO}}_{\mathrm{Fock}_1, \mathrm{Fock}_1}(u) \;=\;
 \frac{(u - h_1)(u - h_2)}{(u + h_1)(u + h_2)}
\]
on charge-$1$ Fock modules.  The UV gluing cocycle is
\[
 \phi^+_{\mathrm{UV}}(u) \;=\;
 \frac{R^{\mathrm{MO}}(u) - \mathrm{id}}{u - u_\star} \, du
 \quad\text{near the wall at } u_\star,
\]
a $1$-form on $\mathfrak{t}^* \simeq \C^2$ valued in $\End(
\mathrm{Fock}_1 \otimes \mathrm{Fock}_1)$, with simple pole at
$u = u_\star$ and residue $R^{\mathrm{MO}}(u_\star)$.  The cocycle
condition on the triple-wall $\{h_1 = 0\} \cap \{h_2 = 0\} \cap \{h_3
= 0\}$ (a point at the origin under the CY$_3$ constraint $h_1 + h_2
+ h_3 = 0$, which reduces the dimension) is the Yang--Baxter equation.

\subsubsection*{The R-matrix on K3 x E}

For $\cC = D^b(\Coh(K3 \times E))$, the spectral parameter is the
elliptic modulus $u = z_1 - z_2 \in E$, and the R-matrix acts on
$\widetilde H^*(K3) \otimes \widetilde H^*(K3) = \C^{4|20} \otimes
\C^{4|20}$.  The MO stable envelope on the derived moduli of
K3-sheaves (Schiffmann--Vasserot \cite{SV2017} for the cohomological
version; Maulik--Toda for the K-theoretic $\hbar$-deformation) gives
an explicit $R$-matrix
\[
 R^{\mathrm{MO}}_{K3}(u) \;=\;
 1 + \hbar \cdot \frac{\Omega_{\mathfrak{osp}(4|20)}}{u} +
 O(\hbar^2)
\]
on $\C^{4|20} \otimes \C^{4|20}$, with $\Omega_{\mathfrak{osp}(4|20)}$
the quadratic Casimir of $\mathfrak{osp}(4|20)$ (the super-Casimir
respecting the Hodge-involution $\theta$).  The spectral parameter
is the elliptic-fibre evaluation-module parameter; the UV gluing
cocycle is
\[
 \phi^+_{\mathrm{UV}, K3 \times E}(u)
 \;=\; \hbar \cdot \frac{\Omega_{\mathfrak{osp}(4|20)}}{u} \, du +
 O(\hbar^2),
\]
a meromorphic $1$-form on $E$ with simple pole at the diagonal
$u = 0$.  The YBE follows by the standard Yangian-super-Yangian
derivation (Arnaudon--Avan--Cramp\'e--Doikou--Frappat--Ragoucy
\cite{AvanFrappat} for $Y_{\mathfrak{osp}}$).

\subsubsection*{The R-matrix on general toric CY_3 with a compact 4-cycle}

For $\cC = \Perf(X)$ with $X$ toric CY$_3$ with a compact $4$-cycle
(LP$^2$ $= \mathrm{Tot}(K_{\P^2})$, $\mathrm{Tot}(K_{\F_n})$, etc.),
the CoHA is given by Rapcak--Soibelman--Yang--Zhao \cite{RSYZ}.  For
LP$^2$, the CoHA is a shuffle algebra on a single generator, and the
Drinfeld double is the quantum toroidal algebra
$U_{q, t}(\widehat{\widehat{\mathfrak{gl}}}_1)$ at a specific
parameter specialisation pinned by the CY$_3$ constraint $h_1 h_2 h_3
\cdot (\text{brane-tensor}) = 0$ (toric fan).  The R-matrix is
Feigin--Jimbo--Miwa's quantum-toroidal $R$-matrix, with spectral
parameter from the $\P^2$ fibre evaluation modules and equivariant
parameters from the CY$_3$ torus action on $X$.


\subsection*{Primary-source table}

Every line of the above derivation has a primary-literature anchor.

\begin{center}
\renewcommand{\arraystretch}{1.3}
\small
\begin{tabular}{rll}
 \toprule
 \textbf{Line} & \textbf{Source} & \textbf{Role} \\
 \midrule
 Dunn--Lurie $E_3 \simeq E_1 \otimes E_2$ &
 Lurie HA 5.1.2.2 & tensor decomposition \\
 Deligne conjecture on $\HH^\bullet$ &
 Kontsevich--Soibelman; Tamarkin; Lurie HA 5.3.1.14 &
 $E_2$-structure on Hochschild \\
 BZFN identification $\HH^0_{\mathrm{cat}} = \cZ$ &
 Ben-Zvi--Francis--Nadler 2010 Thm~4.1 & centre identification \\
 Factorisation homology over $\Sigma_2$ &
 Ayala--Francis 2015 \S 3.1 & Stage-2 specialisation \\
 CoHA = $Y^+(\widehat{\mathfrak{gl}}_1)$ on $\C^3$ &
 Schiffmann--Vasserot 2013 & chiral-side CoHA \\
 Drinfeld double $D(Y^+) = Y(\widehat{\mathfrak{gl}}_1)$ &
 Schiffmann--Vasserot 2013 Thm 8.3 & full Yangian \\
 MO R-matrix for $\mathrm{Hilb}^n$ of surfaces &
 Maulik--Okounkov 2012 Thm 5.1.4 & stable-envelope $R$-matrix \\
 FJMM formula for $R$-matrix on Fock &
 Feigin--Jimbo--Miwa--Mukhin 2017 & explicit Fock $R$ \\
 K3 super-Yangian $Y_{\mathfrak{osp}}(4|20)$ &
 Wave-15 cache + AvFrappat super-Yangian & K3 Drinfeld double \\
 Kontsevich--Tamarkin formality &
 Kontsevich 1999; Tamarkin 1998 & $\mathrm{GRT}_1$ torsor \\
 Costello--Li propagator selects Kontsevich associator &
 Costello--Li 2016 \S 6 & cyclic-sub-torsor pin \\
 CFG 2026 $E_3$ observables &
 Costello--Francis--Gwilliam 2026 & ambient antecedent \\
 Etingof--Kazhdan quantisation of Manin pairs &
 Etingof--Kazhdan I--VI & Drinfeld-double
 construction \\
 Drinfeld New Realisation of Yangian &
 Drinfeld 1986 \emph{Quantum groups}; Kassel GTM~155 &
 universal $R$-matrix \\
 Bezrukavnikov--Finkelberg--Ginzburg Cherednik &
 BFG 2006 \S 6 &
 positive-characteristic Drinfeld double \\
 \bottomrule
\end{tabular}
\end{center}


%%% =====================================================================
\subsection*{Transverse versus longitudinal discipline}
%%% =====================================================================

A summary notational discipline to forbid misreadings.  At $d = 3$,
for a CY$_3$ input $\cC$:

\begin{center}
\renewcommand{\arraystretch}{1.3}
\small
\begin{tabular}{lll}
 \toprule
 \textbf{Axis} & \textbf{Object} & \textbf{Operadic type} \\
 \midrule
 Longitudinal (along $C$) & $A = \Phi_3^{(\Sigma_2, C)}(\cC)$ &
 $E_1$-chiral algebra on $C$ \\
 Transverse (on the centre) &
 $Z^{\mathrm{der}}(\Rep^{E_1}(A))$ &
 $E_2$-algebra (Drinfeld-centre) \\
 Ambient (before Stage-2) & $\cF_\cC = \PhiFA_3(\cC)$ &
 $E_3$-holomorphic factorisation algebra \\
 \bottomrule
\end{tabular}
\end{center}

These are three distinct mathematical objects with three distinct
operadic types.  The transverse $E_2$-shadow is \emph{not} $A$; it
lives on the centre of the representation category of $A$.  The
ambient $E_3$-structure is not $A$ either; $A$ is the
$E_1$-specialisation of the $E_3$-hFA along the curve.  Pattern 273
of the Vol~III ambient-qualifier discipline: the $\Phi$-functor has
three object-level outputs depending on scope, the transverse one
being the $E_2$-shadow documented in this section.

This closes the chain-level derivation of the transverse $E_2$-shadow.

%%% =====================================================================
\subsection*{Appendix A: detailed Deligne-Kontsevich brace formulas}
%%% =====================================================================

For explicitness we record the brace-operation formulas giving the
$E_2$-action on Hochschild cochains.  For $\phi \in \mathrm{HH}^p(A,
A)$ and $\psi \in \mathrm{HH}^q(A, A)$, the brace operation $\phi\{\psi\}$
is the degree-$(p+q-1)$ cochain
\[
 (\phi\{\psi\})(a_1, \ldots, a_{p+q-1}) \;=\;
 \sum_{i=1}^{p} (-1)^{(q-1)\epsilon_i}
 \phi(a_1, \ldots, a_{i-1}, \psi(a_i, \ldots, a_{i+q-1}),
 a_{i+q}, \ldots, a_{p+q-1}),
\]
with sign $\epsilon_i$ governed by the Koszul convention.  The
Gerstenhaber bracket is
\[
 [\phi, \psi] \;=\; \phi\{\psi\} - (-1)^{(p-1)(q-1)} \psi\{\phi\},
\]
which, at the cohomology level, is the Lie bracket on
$\mathrm{HH}^\bullet(A, A)$.  The cup product is
$(\phi \smile \psi)(a_1, \ldots, a_{p+q}) = (-1)^{pq} \phi(a_1, \ldots,
a_p) \cdot \psi(a_{p+1}, \ldots, a_{p+q})$.

McClure--Smith \cite{McClureSmith} upgrade the pair
$(\text{cup product}, \text{brace})$ to an $E_2$-operad action on
$\CC^\bullet(A, A)$; Kontsevich--Soibelman
\cite{KontsevichSoibelmanDefI} and Tamarkin \cite{TamarkinDGCat}
upgrade to the full homotopy-$E_2$ coherences.

For an $E_1$-chiral algebra $A$ along $C$, the ``multiplication''
$\mu \colon A \otimes A \to A$ is replaced by the OPE-bracket
$\mu^{\mathrm{ch}} \colon j_* j^*(A \boxtimes A) \to \Delta_* A$ (BD
chiral-algebra datum; Beilinson--Drinfeld \cite{BeilinsonDrinfeldCA}),
and the Hochschild complex $\mathrm{HH}^\bullet(A, A)$ is replaced by
the chiral-Hochschild complex
$\mathrm{ChirHoch}^\bullet(A, A)$.  The Deligne conjecture in this
setting is Francis \cite[\S 2]{Francis2013}: the chiral-Hochschild
complex carries a canonical $E_2$-action extending the brace
operations.

%%% =====================================================================
\subsection*{Appendix B: Cross-verification against Vol III cache entries}
%%% =====================================================================

The chain-level derivation above is in direct concordance with the
following Vol~III cache entries and the broader Vol~III knowledge
graph.  Each entry corresponds to a residual obstruction or
disambiguation that the above treatment respects.

\begin{itemize}
\item \textbf{Cache entry: ChirHoch vs THH (structural, not dimensional).}
Cache memory entries 46--50.  We distinguished
$\mathrm{ChirHoch}(A, A)$ (chiral-Hochschild, BD-chiral derived centre)
from $\mathrm{THH}(A)$ (topological-Hochschild, $\mathrm{SC}^{\mathrm{ch,
top}}$ derived centre); the first is the transverse $E_2$-shadow of the
longitudinal $E_1$-chiral algebra $A$, the second is the $\mathrm{SC}^{
\mathrm{ch, top}}$-analogue relevant to Vol~II.

\item \textbf{Cache entry: BZFN different algebras.}  The BZFN
equivalence $\mathrm{HH}^\bullet_{\mathrm{cat}}(\cC, \cC) \simeq
\cZ^{\mathrm{der}}(\cC)$ of dg-categories is applied with $\cC =
\Rep^{E_1}(A)$ (not $\cC = \Perf(A)$); this matches the cache memory
entry ``BZFN uses different algebras, not different ambient
categories.''

\item \textbf{Cache entry: Spectral parameter from evaluation modules.}
$R^{\mathrm{ch}}(u)$ acts on evaluation modules $V_{u_1} \otimes
V_{u_2}$; the spectral parameter is $u = u_1 - u_2$, not a variable of
the Drinfeld centre itself.  The Drinfeld centre has no spectral
parameter; the universal $R$-matrix is in
$Z^{\mathrm{der}} \widehat\otimes Z^{\mathrm{der}}$, and its evaluation
on a pair of evaluation modules produces the spectral-dependent $R$.

\item \textbf{Cache entry: CoHA $=$ $Y^+$, double $=$ $Y$.}  We
consistently distinguished $\mathrm{CoHA}(\C^3) = Y^+(\widehat{\mathfrak
{gl}}_1)$ (associative, half) from $Y(\widehat{\mathfrak{gl}}_1)$
(Hopf, full double); the Drinfeld-centre equivalence is with the
double, not with the half.

\item \textbf{Cache entry: K3 super-Yangian $Y_{\mathfrak{osp}}(4|20)$.}
The signature $(4, 20)$ is forced by the Mukai Hodge involution on
$\widetilde H^*(K3)$, matching the Wave-15 project memory
\texttt{project\_super\_yangian\_Y\_osp\_4\_20.md}.

\item \textbf{Cache entry: two-stage factorisation + cyclic
$\mathrm{GRT}_1$.}  The ambiguity in the Dunn--Lurie decomposition is
the cyclic $\mathrm{GRT}_1$-torsor, pinned to the Kontsevich associator
by the Costello--Li propagator (Wave-12 memory
\texttt{project\_two\_stage\_factorisation.md}).

\item \textbf{Cache entry: four $\kappa$ discipline.}  The transverse
$E_2$-shadow carries no $\kappa$-invariant of its own; the ambient
$\kappa$-invariants are:
$\kappa_{\mathrm{ch}}(A)$ on the longitudinal $E_1$-algebra;
$\kappa_{\mathrm{cat}}(\cC)$ on the CY$_3$-category; $\kappa_{\mathrm{BKM}}$
on the BKM side; $\kappa_{\mathrm{fiber}}$ for lattice corrections.
These are not invariants of the transverse shadow itself; they are
invariants of adjacent structures that feed into the transverse
shadow through the Stage-2 specialisation.

\item \textbf{Cache entry: Pattern 273 (Phi as functor vs object-level).}
The transverse $E_2$-shadow is an object-level construction; it does
not promote to an $(\infty, 1)$-functor from CY$_3$ to $E_2$-algebras
without further morphism-preservation data.  At $(\infty, 1)$-level,
the corresponding functor is a scope declaration; object-level
correctness is what we establish above.

\item \textbf{Wave-14 insight: hCS geometric resolution.}  The
holomorphic Chern--Simons theory is a geometric resolution of the
categorical $E_3$-structure on $\HH^\bullet(\cC, \cC)$; our chain-level
transverse $E_2$-shadow is the $E_2$-component of this $E_3$-structure
after Dunn--Lurie.  This aligns with Wave-14 memory
\texttt{project\_hcs\_vs\_cat\_hochschild\_identification.md}.
\end{itemize}

%%% =====================================================================
\subsection*{Appendix C: Status within the five-theorem skeleton}
%%% =====================================================================

The transverse $E_2$-shadow touches three of the five Vol~III theorems.

\emph{Theorem B (chiral Positselski, Koszul duality).}  The
$E_2$-enhanced Koszul duality on the transverse shadow is the
categorical statement of CY-B$_3$:
$\cZ(\Rep^{E_1}(A)) \simeq \cZ(\Rep^{E_1}(A^!))^{\mathrm{rev}}$ as
braided monoidal categories (Theorem~\ref{thm:cy-b-d3}).  The
transverse $E_2$-shadow is where this equivalence lives.

\emph{Theorem C (derived-centre complementarity).}  The
$\kappa_{\mathrm{ch}} + \kappa^!_{\mathrm{ch}} \in \{0, 8, 13, 250/3,
98/3\}$ landmark ceiling is on the longitudinal $A$ and $A^!$, not on
the transverse shadow.  The transverse shadow provides the
braided-tensor environment in which $A^!$ is constructed.

\emph{Theorem H (Hochschild concentration).}  The Hochschild cochain
complex $\mathrm{HH}^\bullet(A, A)$ of the longitudinal algebra $A$
is precisely the $E_2$-algebra of the transverse shadow; the
concentration result of Vol~I Theorem~H is the statement that this
complex is concentrated in finitely many degrees (for suitable $A$).

At the $(\infty, 1)$-level, the transverse shadow is a target of the
$\Phi$-functor-as-$(\infty, 1)$-functor declaration in Pattern 273:
$\Phi_{E_2}(\cC) = Z^{\mathrm{ch}}(\Phi_{E_1}(\cC))$ is precisely the
conjectural statement that the transverse shadow assignment lifts to
a functorial construction on the CY$_3$-category of inputs
(Conjecture~\ref{conj:phi-e2-drinfeld-center}).  The chain-level
content above is the object-level case of this conjecture, verified
case-by-case for the three verified CY$_3$ inputs.
